\chapter{References}
\label{ch:references}
\renewcommand{\bibname}{}
\begingroup
\renewcommand{\chapter}[2]{}% the chapter heading above replaces thebibliography's own
\small
\begin{thebibliography}{99}
\setlength{\itemsep}{3pt}

\bibitem{zhang2022buildsheriff} C. Zhang, B. Chen, X. Peng, and W. Zhao, ``BuildSheriff: Change-aware test failure triage for continuous integration builds,'' in \emph{Proc. 44th Int. Conf. Software Engineering (ICSE)}, 2022. [Online]. Available: \url{https://chenbihuan.github.io/paper/icse22-zhang-buildsheriff.pdf}

\bibitem{saidani2020predicting} I. Saidani, A. Ouni, M. Chouchen, and M. W. Mkaouer, ``Predicting continuous integration build failures using evolutionary search,'' \emph{Information and Software Technology}, 2020. [Online]. Available: \url{https://ouniali.github.io/papers/IST2020-saidani.pdf}

\bibitem{barrak2021why} A. Barrak, E. E. Eghan, B. Adams, and F. Khomh, ``Why do builds fail? A conceptual replication study,'' \emph{Journal of Systems and Software}, 2021. [Online]. Available: \url{https://www.sciencedirect.com/science/article/abs/pii/S0164121221000364}

\bibitem{rausch2017empirical} T. Rausch, W. Hummer, P. Leitner, and S. Schulte, ``An empirical analysis of build failures in the continuous integration workflows of Java-based open-source software,'' in \emph{Proc. 14th Int. Conf. Mining Software Repositories (MSR)}, 2017, pp.~345--355. [Online]. Available: \url{https://ieeexplore.ieee.org/document/7962384}

\bibitem{durieux2020empirical} T. Durieux, C. Le Goues, M. Hilton, and R. Abreu, ``Empirical study of restarted and flaky builds on Travis CI,'' in \emph{Proc. 17th Int. Conf. Mining Software Repositories (MSR)}, 2020. [Online]. Available: \url{https://2020.msrconf.org/details/msr-2020-papers/21}

\bibitem{ghaleb2019empirical} T. A. Ghaleb, D. A. da Costa, and Y. Zou, ``An empirical study of the long duration of continuous integration builds,'' \emph{Empirical Software Engineering}, 2019. [Online]. Available: \url{https://link.springer.com/article/10.1007/s10664-019-09695-9}

\bibitem{saidani2022improving} I. Saidani, A. Ouni, and M. W. Mkaouer, ``Improving the prediction of continuous integration build failures using deep learning,'' \emph{Automated Software Engineering}, 2022. [Online]. Available: \url{https://www.researchgate.net/publication/357990835}

\bibitem{farihane2025ci} R. Farihane, ``CI/CD pipeline optimization using AI: A systematic mapping study,'' \emph{MDPI Proceedings}, 2025.

\bibitem{kumar2025optimizing} R. Kumar, D. N. D. Kollipara, A. H. Mohammed, S. Kumar, F. B. Mohammed, and A. Kumar, ``Optimizing CI/CD pipelines with AI-driven build failure prediction: An empirical study on machine learning models for early failure detection,'' 2025.

\bibitem{hong2024practitioners} Y. Hong, C. Tantithamthavorn, J. Pasuksmit, P. Thongtanunam, A. Friedman, X. Zhao, and A. Krasikov, ``Practitioners' challenges and perceptions of CI build failure predictions at Atlassian,'' arXiv preprint arXiv:2402.09651, 2024.

\bibitem{hilton2016usage} M. Hilton, T. Tunnell, K. Huang, D. Marinov, and D. Dig, ``Usage, costs, and benefits of continuous integration in open-source projects,'' in \emph{Proc. 31st IEEE/ACM Int. Conf. Automated Software Engineering (ASE)}, 2016, pp.~426--437.

\bibitem{beller2017travistorrent} M. Beller, G. Gousios, and A. Zaidman, ``TravisTorrent: Synthesizing Travis CI and GitHub for full-stack research on continuous integration,'' in \emph{Proc. 14th Int. Conf. Mining Software Repositories (MSR)}, 2017, pp.~447--450.

\bibitem{hassan2017change} F. Hassan and X. Wang, ``Change-aware build prediction model for stall avoidance in continuous integration,'' in \emph{Proc. ACM/IEEE Int. Symp. Empirical Software Engineering and Measurement (ESEM)}, 2017, pp.~157--162.

\bibitem{ni2017cost} A. Ni and M. Li, ``Cost-effective build outcome prediction using cascaded classifiers,'' in \emph{Proc. 14th Int. Conf. Mining Software Repositories (MSR)}, 2017, pp.~455--458.

\bibitem{xia2017could} J. Xia and Y. Li, ``Could we predict the result of a continuous integration build? An empirical study,'' in \emph{Proc. IEEE Int. Conf. Software Quality, Reliability and Security Companion (QRS-C)}, 2017, pp.~311--315.

\bibitem{chen2020buildfast} B. Chen, L. Chen, C. Zhang, and X. Peng, ``BuildFast: History-aware build outcome prediction for fast feedback and reduced cost in continuous integration,'' in \emph{Proc. 35th IEEE/ACM Int. Conf. Automated Software Engineering (ASE)}, 2020, pp.~42--53.

\bibitem{jin2020cost} X. Jin and F. Servant, ``A cost-efficient approach to building in continuous integration,'' in \emph{Proc. 42nd ACM/IEEE Int. Conf. Software Engineering (ICSE)}, 2020, pp.~13--25.

\bibitem{gallaba2018noise} K. Gallaba, C. Macho, M. Pinzger, and S. McIntosh, ``Noise and heterogeneity in historical build data: An empirical study of Travis CI,'' in \emph{Proc. 33rd ACM/IEEE Int. Conf. Automated Software Engineering (ASE)}, 2018, pp.~87--97.

\bibitem{durieux2019analysis} T. Durieux, R. Abreu, M. Monperrus, T. F. Bissyandé, and L. Cruz, ``An analysis of 35+ million jobs of Travis CI,'' in \emph{Proc. IEEE Int. Conf. Software Maintenance and Evolution (ICSME)}, 2019, pp.~291--295.

\bibitem{kaufman2012leakage} S. Kaufman, S. Rosset, C. Perlich, and O. Stitelman, ``Leakage in data mining: Formulation, detection, and avoidance,'' \emph{ACM Trans. Knowledge Discovery from Data}, vol.~6, no.~4, 2012.

\bibitem{bergmeir2012use} C. Bergmeir and J. M. Benítez, ``On the use of cross-validation for time series predictor evaluation,'' \emph{Information Sciences}, vol.~191, pp.~192--213, 2012.

\bibitem{nosek2018preregistration} B. A. Nosek, C. R. Ebersole, A. C. DeHaven, and D. T. Mellor, ``The preregistration revolution,'' \emph{Proc. National Academy of Sciences}, vol.~115, no.~11, pp.~2600--2606, 2018.

\bibitem{breiman2001random} L. Breiman, ``Random forests,'' \emph{Machine Learning}, vol.~45, no.~1, pp.~5--32, 2001.

\bibitem{quinlan1986induction} J. R. Quinlan, ``Induction of decision trees,'' \emph{Machine Learning}, vol.~1, no.~1, pp.~81--106, 1986.

\bibitem{cortes1995support} C. Cortes and V. Vapnik, ``Support-vector networks,'' \emph{Machine Learning}, vol.~20, no.~3, pp.~273--297, 1995.

\bibitem{chen2016xgboost} T. Chen and C. Guestrin, ``XGBoost: A scalable tree boosting system,'' in \emph{Proc. 22nd ACM SIGKDD Int. Conf. Knowledge Discovery and Data Mining}, 2016, pp.~785--794.

\bibitem{ke2017lightgbm} G. Ke, Q. Meng, T. Finley, T. Wang, W. Chen, W. Ma, Q. Ye, and T.-Y. Liu, ``LightGBM: A highly efficient gradient boosting decision tree,'' in \emph{Advances in Neural Information Processing Systems (NeurIPS)}, vol.~30, 2017.

\bibitem{bruch2021alternative} S. Bruch, ``An alternative cross entropy loss for learning-to-rank,'' in \emph{Proc. Web Conference (WWW)}, 2021, pp.~118--126.

\bibitem{platt1999probabilistic} J. C. Platt, ``Probabilistic outputs for support vector machines and comparisons to regularized likelihood methods,'' in \emph{Advances in Large Margin Classifiers}. MIT Press, 1999, pp.~61--74.

\bibitem{lundberg2017unified} S. M. Lundberg and S.-I. Lee, ``A unified approach to interpreting model predictions,'' in \emph{Advances in Neural Information Processing Systems (NeurIPS)}, vol.~30, 2017.

\bibitem{liu2008isolation} F. T. Liu, K. M. Ting, and Z.-H. Zhou, ``Isolation forest,'' in \emph{Proc. 8th IEEE Int. Conf. Data Mining (ICDM)}, 2008, pp.~413--422.

\bibitem{pedregosa2011scikit} F. Pedregosa \emph{et al.}, ``Scikit-learn: Machine learning in Python,'' \emph{Journal of Machine Learning Research}, vol.~12, pp.~2825--2830, 2011.

\bibitem{salton1988term} G. Salton and C. Buckley, ``Term-weighting approaches in automatic text retrieval,'' \emph{Information Processing \& Management}, vol.~24, no.~5, pp.~513--523, 1988.

\bibitem{lloyd1982least} S. P. Lloyd, ``Least squares quantization in PCM,'' \emph{IEEE Trans. Information Theory}, vol.~28, no.~2, pp.~129--137, 1982.

\bibitem{efron1993introduction} B. Efron and R. J. Tibshirani, \emph{An Introduction to the Bootstrap}. Chapman \& Hall, 1993.

\end{thebibliography}
\endgroup

\vspace{6pt}
\noindent\textbf{Project resources.} Source code, protocols and results: \url{https://github.com/yasaswinipakam/cicd-ai-optimizer}. Read-only demonstration: \url{https://cicd-ai-optimizer.onrender.com}.
